\section{One-sided limits}

A two-sided limit asks what happens as \(x\) approaches a point from both sides. Sometimes the behavior from the left and from the right is different. To describe this, we use one-sided limits.

The left-hand limit is written as
\[
\lim_{x\to a^-} f(x).
\]
It describes what happens when \(x\) approaches \(a\) through values smaller than \(a\).

The right-hand limit is written as
\[
\lim_{x\to a^+} f(x).
\]
It describes what happens when \(x\) approaches \(a\) through values greater than \(a\).

\subsection*{When the two-sided limit exists}

A two-sided limit exists exactly when the left-hand and right-hand limits both exist and are equal.

\begin{theorembox}
The limit
\[
\lim_{x\to a}f(x)=L
\]
exists if and only if
\[
\lim_{x\to a^-}f(x)=L
\qquad\text{and}\qquad
\lim_{x\to a^+}f(x)=L.
\]
\end{theorembox}

If the two one-sided limits are different, then the two-sided limit does not exist.

\begin{examplebox}
Consider
\[
f(x)=\frac{|x|}{x},
\qquad x\ne0.
\]
If \(x>0\), then \(|x|=x\), so
\[
f(x)=1.
\]
If \(x<0\), then \(|x|=-x\), so
\[
f(x)=-1.
\]
Therefore
\[
\lim_{x\to0^-}\frac{|x|}{x}=-1,
\qquad
\lim_{x\to0^+}\frac{|x|}{x}=1.
\]
The one-sided limits are different, so
\[
\lim_{x\to0}\frac{|x|}{x}
\]
does not exist.
\end{examplebox}

This example shows why looking from both sides matters. A formula may look simple, but its behavior can depend on the direction of approach.

\subsection*{Piecewise functions}

One-sided limits are especially important for piecewise-defined functions. Each side of the point may use a different formula.

\begin{examplebox}
Let
\[
f(x)=
\begin{cases}
x+2, & x<1,\\
3x, & x\ge1.
\end{cases}
\]
The left-hand limit at \(1\) uses the formula \(x+2\):
\[
\lim_{x\to1^-}f(x)=1+2=3.
\]
The right-hand limit at \(1\) uses the formula \(3x\):
\[
\lim_{x\to1^+}f(x)=3\cdot1=3.
\]
The two one-sided limits are equal, so
\[
\lim_{x\to1}f(x)=3.
\]
Also, because \(f(1)=3\cdot1=3\), the function is continuous at \(1\). We will define continuity later in this chapter.
\end{examplebox}

\begin{examplebox}
Let
\[
g(x)=
\begin{cases}
x+2, & x<1,\\
3x+1, & x\ge1.
\end{cases}
\]
Then
\[
\lim_{x\to1^-}g(x)=3,
\]
but
\[
\lim_{x\to1^+}g(x)=4.
\]
Since the one-sided limits are different,
\[
\lim_{x\to1}g(x)
\]
does not exist.
\end{examplebox}

The point of the calculation is not only to obtain numbers. The calculation tells us whether the two sides of the graph meet at the same height.

\subsection*{Reading from a graph}

On a graph, the left-hand limit is the height approached as we move toward the point from the left. The right-hand limit is the height approached as we move toward the point from the right.

If the graph approaches the same height from both sides, the two-sided limit exists. If the graph approaches different heights, there is a jump.

\begin{summarybox}
When studying one-sided limits, ask:
\begin{itemize}
\item What happens when \(x\) approaches from the left?
\item What happens when \(x\) approaches from the right?
\item Are the two one-sided limits equal?
\item Does the formula change on different sides of the point?
\item Does the graph show a jump or a meeting point?
\end{itemize}
\end{summarybox}

One-sided limits are the language we need when behavior depends on direction.